\section*{Executive Summary}
\addcontentsline{toc}{section}{Executive Summary}

The portfolio consultation evaluates a diversified portfolio consisting of three Philippine equities: PLDT Inc. (TEL), Manila Electric Company (MER), and Aboitiz Equity Ventures (AEV). In addition, it includes international technology equities NVIDIA Corporation (NVDA) and Meta Platforms, Inc. (META) as well as Bitcoin (BTC). Using daily market data from January 2020 through December 2025, the study assesses portfolio performance, diversification, optimization, and downside risk. Furthermore, the SPDR S\&P 500 ETF Trust (SPY) serves as the benchmark while the Philippine 91-day Treasury Bill serves as the risk-free rate.

The empirical results demonstrate that the portfolio substantially outperformed the benchmark. However, this out-performance came with significantly greater risk. The equal-weighted Buy-and-Hold strategy generated a 743.04\% cumulative return and a 42.83\% CAGR. In comparison, SPY generated a 129.06\% cumulative return and a 14.87\% CAGR. Nevertheless, the Buy-and-Hold strategy recorded 35.12\% annualized volatility and a -58.21\% maximum drawdown. Meanwhile, monthly rebalancing reduced annualized volatility to 24.18\% and maximum drawdown to -43.20\%. However, cumulative returns also declined to 384.85\%.

Moreover, the downside-risk analysis highlights the trade-off between return maximization and capital preservation. During major market stress events, the GMV portfolio consistently experienced smaller losses than the Maximum Sharpe portfolio. For instance, during the 2020 COVID-19 crash, GMV declined by 22.55\% compared with 28.83\% for Maximum Sharpe. Similarly under the hypothetical 1.5x Black Swan scenario, losses were 35.89\% and 45.30\%, respectively.

Therefore, portfolio selection should reflect the investor's risk tolerance and objectives. The Maximum Sharpe portfolio suits investors seeking higher returns despite greater risk. Conversely, the GMV portfolio suits risk-averse investors prioritizing capital preservation. For a balance between growth and risk control, monthly rebalancing provides a disciplined approach. It reduces volatility and drawdown while still substantially outperforming the benchmark.